\section{Returning to the Research Questions}
The evidence answers the research questions at different levels of certainty. RQ1 asked whether a correct first-minute dementia label is necessarily stable across the full examined interval. The original SVM results show that it is not: reported first-minute balanced accuracy was 73.48\% for AD versus controls and 73.72\% for FTD versus controls, while substantial fractions of predictions changed during the five-minute observation. Stable but wrong decisions also occurred. These are descriptive historical results in a 64-person analytic cohort and require complete source-table reconciliation before an independent replication claim.

RQ2 asked whether the custom DSRF decoder with a frozen CBraMod backbone is competitive. The answer is yes in a limited internal sense: DSRF obtained 77.06\% pooled development BAcc and 40 correct decisions in 52 participants, almost identical to XGBoost's 76.61\% and 40 correct decisions. However, the final locked 13-person assessment favored XGBoost, 77.38\% versus 70.24\%. We cannot call the decoder superior on the strength of the development result alone.

RQ3 asked what additional reliability tests reveal. Artificial Gaussian noise reduced DSRF's BAcc and changed its predictions; the channel-zeroing test caused fewer changes but does not establish robustness to real missing electrodes. Input masking found channel-specific sensitivities that cannot be interpreted as biological causes. The separate MCI experiment showed why stable but wrong predictions are a necessary negative control. RQ4 concerned generalization: small cohorts, uncertain overlap between study tracks and the lack of external AD validation leave important limits unresolved.

\section{Why the Two-Part Design is Coherent}
At first glance a retrospective SVM stability study and a pretrained foundation-model decoder could look like unrelated projects. They answer two parts of one practical research problem. Before asking whether a new model is better, we need to know what ``better'' should mean. A single classification score can hide whether a model changes its answer under a different part of the recording or remains confident when it is incorrect. This was the purpose of the earlier temporal work.

The second part asks whether a more flexible feature representation can be useful when evaluated according to a stricter standard. CBraMod provides pretrained EEG representations. DSRF adds task-specific temporal features and a gated residual pathway. Subject-wise development, a locked final test, baseline comparisons and artificial stress tests then establish the decoder's capabilities and weaknesses. Thus, the decoder is an architectural investigation motivated by the reliability problem, not a claim that the earlier SVM experiment directly trained or validated EEGraft.

This also explains the title's grammar. ``A Gated Dual-Scale Residual Fusion Decoder \emph{and} Subject-Level Reliability Assessment'' describes two connected contributions. It does not state that a single trained decoder passed every stability criterion. We use the EEGraft name for the decoder-centered research approach while keeping the original spectral analysis clearly identified.

\section{Why XGBoost Matters}
The strongest classical comparator reached a higher BAcc, accuracy, macro-F1 and AUROC on the locked final group than the proposed decoder. Because XGBoost is not a foundation-model architecture, this outcome places a real boundary on the interpretation of representation learning. Feature-based models may remain competitive on a small, relatively uniform EEG dataset. A more complex model can be harder to optimize or can adapt to development samples in ways that do not improve a specific unseen group. The present data do not isolate which explanation is correct. We therefore cannot ascribe the final difference to overfitting, architecture design or pretraining mismatch without further controlled experiments.

The paired interval of $-28.57$ to $+14.29$ percentage points for DSRF minus XGBoost also matters. It spans both meaningful disadvantages and possible small advantages. In other words, the observed XGBoost lead is a fact about these 13 people, but the true relative performance on an independent population is not known precisely. Varoquaux's discussion of small-sample error bars is particularly relevant here \cite{varoquaux2018}. This finding is not a failure of honest reporting; it is the appropriate limit of the available evidence.

\section{Model Complexity and Interpretability}
DSRF learns an interaction between pretrained high-level tokens and lower-level temporal filters. The gate can adapt the contribution of one path, and the residual retains a direct information route. The low-only model's weaker development result is consistent with useful information in pretrained representations, while the matched blend comparison suggests the form of the residual may affect behavior. These statements are comparative observations rather than proof that the mathematical components represent distinct neural physiology.

Interpretability has a similar boundary. Electrode masking may show which input changes disturb the model's output, but there is no one-to-one mapping between sensitivity and neuronal disease mechanism. The channel-occlusion explanation is best used to ask whether the trained system depends excessively on a few electrodes or acquisition patterns. It should not be used to claim discovery of disease-specific brain biomarkers. Future tests can investigate whether sensitivity profiles remain stable across individuals, hardware and hospitals.

\section{Internal Validity}
Internal validity asks whether the stated protocol supports the reported comparison within the analyzed data. Important strengths include participant-disjoint splits, use of balanced accuracy for class-imbalanced comparisons, a frozen final holdout and explicit reporting of negative results. The final protocol records pretrained checkpoint hashes and exact model settings, improving traceability. The paired holdout results are sufficiently clear to reconstruct each model's confusion matrix.

Potential weaknesses remain. Dataset transformations may contain preprocessing choices that require inspection in the actual executed code rather than summary text alone. The early 64-person temporal project was summarized retrospectively, and its original machine-readable tables and selected participants require reconciliation. Development-fold model selection and exploratory ablation choices introduce selection bias. Cross-validation fold scores are correlated because the same source cohort contributes to multiple comparisons. None of these limitations should be ignored simply because the final test was not directly used for hyperparameter tuning.

\section{External Validity}
External validity is the weakest part of the evidence. The later DSRF and XGBoost comparisons draw from a single public source cohort. Their participants were recorded under a limited set of clinical and hardware conditions. They do not constitute a cross-hospital trial or an independent test on a separate AD dataset. The Figshare compatibility audit could not establish a clean 19-channel AD-versus-control sample under the same protocol. The Isfahan recording analysis was an MCI-versus-control task with additional channel-unit and participant-identity uncertainty. Neither provides external validation of EEGraft for AD.

The 2026 benchmark literature emphasizes the practical importance of external distributions, differing montages and transparent participant-level evaluation \cite{miltiadous2026benchmark,masry2026,xiong2026}. A future extension should not only obtain another file collection; it must establish the number of distinct people, harmonize channel metadata and sampling, apply the same model without tuning on the new labels, and report failures as well as successes.

\section{Statistical and Measurement Limitations}
Balanced accuracy is intuitive but unstable with only seven AD and six control cases in the heldout test. Macro-F1 is also sensitive to individual class assignments. AUROC uses probabilities, but with 13 people its value can change substantially if one participant moves in rank. The Brier score and ECE add useful context but cannot prove calibrated clinical risk in such a small sample. The confidence intervals are exploratory within-cohort bootstrap intervals; they do not automatically represent uncertainty due to all possible model training seeds, dataset shifts or prospective sampling.

The temporal change rate also needs care. It is affected by the number of windows, how long the recording lasts and the threshold used to assign classes. A model that predicts one class for every person has perfect stability but only 50\% balanced accuracy in a two-class task. A model with no changes should therefore not receive a reliability award unless it also demonstrates discrimination. The MCI negative controls made this conceptual problem visible in real saved experiments.

\section{Practical Meaning and Public Health Claims}
The thesis demonstrates why EEG classification should be examined beyond a single score. It does not demonstrate a deployable screening tool, a treatment, a prevention technique or a forecast of future Alzheimer's disease. A positive or negative classification in these retrospective recordings cannot tell an asymptomatic person how likely they are to develop AD years later. Such a claim would require longitudinal participants, explicit future outcomes, calibration against disease incidence and clinical validation. The current labels reflect already assigned diagnostic categories.

The modest and uncertain test performance also means the method should not be used for patient decisions. A clinician needs information about symptoms, medical history and established assessments that these EEG-only experiments did not test. The most useful immediate audience is researchers choosing evaluation protocols and architectures for future studies.

\section{Lessons for Reproducibility and Presentation}
A thesis can be technically accurate yet confusing if it hides changes in cohorts or calls all classification experiments ``the model.'' We avoid that ambiguity through explicit Study A and Study B labels, a cohort table, and separate result sections. The figures are accompanied by captions that state what is observed and what remains illustrative. The bibliography includes both architecture papers and validation research, because methodology is central to this topic.

Some source documents and previously approved figure files were not accessible in the current handoff. The report package therefore contains available user-provided illustrations and reproducible graphics made from exact saved aggregate numbers, with provenance files. Those graphics should not be mistaken for recovered original raw EEG traces. Before submission, a final audit should verify image labels, source tables and historical model settings against the original experiment directory. This is a document-completion issue, not a reason to fabricate replacement evidence.

\section{Overall Interpretation}
The strongest supported conclusion is about measurement. A model's correctness at one interval, its consistency across time, its sensitivity to artificial corruption and its calibration are distinct properties. A custom decoder can be an interesting engineering contribution even if it does not surpass a simpler baseline on the final test. The project is most defensible when it reports both insights together and leaves untested clinical claims out of the conclusion.
